\section{Empirical Evaluation}

\subsection{Performance Metrics}
We evaluate the algorithms based on the following metrics:

\begin{itemize}
    \item \textbf{Caption quality}: At the testing phase, we evaluate the generated difference captions using standard metrics, including BLEU-4, METEOR, ROUGE-L, and CIDEr-D. These metrics measure how well the generated descriptions align with human annotations.
\end{itemize}

\subsection{Baseline Methods}
We consider the following recent state-of-the-art method as our primary baseline:

\begin{itemize}
    \item \textbf{ADPO}: An adversarial direct preference optimization framework that formulates image difference captioning (IDC) as a preference learning problem. It improves fine-grained difference alignment by coupling the captioning model with an adversarially trained hard-negative retriever to generate challenging counterfactual samples.
\end{itemize}

ADPO is a strong baseline in IDC and represents a competitive preference optimization framework for vision-language alignment.

\subsection{Benchmark Tasks or Datasets}
We evaluate the algorithms on several benchmark datasets widely used in the image difference captioning (IDC) literature:

\begin{itemize}
    \item \textbf{CLEVR-Change}: A synthetic dataset for scene-level difference captioning, where the model must identify semantic changes (e.g., object color, shape, and position) while ignoring non-semantic variations such as viewpoint and illumination changes.

    \item \textbf{Spot-the-Diff}: A real-world image-pair dataset that requires the model to generate natural language descriptions of semantic differences between two images.

    \item \textbf{Image-Editing-Request}: A dataset consisting of image pairs and corresponding editing instructions, where the model needs to describe fine-grained and diverse editing operations between images.
\end{itemize}